% RESEARCH_DRAFT_SHA256: 585859db6e82107b086b50f2ef8562581ed259bb4dc2f2d0ecf8885f3eaf4be8
% Layout and source-reference wrappers expanded for the web renderer.
\documentclass[11pt]{article}
\usepackage[margin=25mm]{geometry}
\usepackage{amsthm}
\usepackage{fontspec,unicode-math}
\setmainfont{Latin Modern Roman}
\setmathfont{Latin Modern Math}
\usepackage{amsmath,amssymb,mathtools}
\usepackage{xurl,hyperref}
\hypersetup{colorlinks=true,urlcolor=blue,linkcolor=blue}
\setcounter{secnumdepth}{0}
\setlength{\parindent}{0pt}
\setlength{\parskip}{5pt}
\setlength{\emergencystretch}{3em}
\newcommand{\sref}[2]{\hyperlink{source:#1:#2}{[S#2]}}
\newcommand{\cataloguescope}[1]{\paragraph{Source question.}#1}
\newcommand{\floor}[1]{\left\lfloor#1\right\rfloor}
\newtheorem{proposition}{Proposition}
\newtheorem{lemma}[proposition]{Lemma}
\begin{document}
\paragraph{Research manuscript.} OpenAI \texttt{gpt-6-astra}, reasoning \texttt{ultra}; 11 October 2026.
\paragraph{Preserved source body.} The complete catalogue mathematical text and sources follow with layout and source links expanded, including source qualifications and its historical preparation notes. New research begins at the section headed Research attempt.

\section{Geometric Explanations and Polynomial Pivot Rules for Simplex}\label{statement}

\subsection{Definitions and mathematical statement}

\paragraph{Shared notation.}Standard number systems are denoted by $\mathbb N,\mathbb Z,\mathbb Q,\mathbb R,\mathbb C$; index ranges, computational models and specialized notation follow the cited problem source and the definitions below.


Polynomial time means a number of elementary steps bounded by a polynomial in the encoded input length. A decision problem is in $\mathrm{NP}$ if yes-instances have polynomial-length certificates verifiable in polynomial time; NP-hardness is under polynomial-time reductions.

A linear program optimizes $c^Tx$ over a polyhedron $\{x:Ax\leq b\}$. In the simplex method a pivot moves between adjacent basic feasible solutions; a pivot rule chooses the next such move.
\paragraph{Statement recorded in the supplied source.}
For the simplex method for linear programming, can one explain its good practical performance by a geometric property of real-world linear programs, and does there exist a polynomial-time pivot rule? \href{https://doi.org/10.4171/owr/2018/54}{[S2]}
\paragraph{Scope qualification.}
The practical-explanation subquestion supplies no dataset or formal criterion. The pivot-rule question retains the source\textquotesingle s permitted rule and arithmetic/bit-complexity conventions.
\paragraph{Source status.}
The supplied catalogue labels this record open. Its source-status metadata was checked against the supplied snapshot on 2026-10-11. This import does not independently establish current open status or mathematical correctness. The catalogue literature assessment is dated 2026-08-21.
\subsection{Short English statement}
Explain the cited practical simplex behavior geometrically and find a polynomial-time pivot rule.
\subsection{Sources}
\begin{itemize}
\item[S1] Ulam AI and the UnsolvedMath Contributors, UnsolvedMath: A Curated Collection of Open Mathematics Problems, catalogue entry. Catalogue curation and metadata: CC BY 4.0. The supplied source question is retained; definitions and the English summary are editorial additions.
\url{https://huggingface.co/datasets/ulamai/UnsolvedMath}.
\item[S2] Convex Geometry and its Applications (2019), Oberwolfach Reports, DOI 10.4171/owr/2018/54, pp. 3217--3281
\url{https://doi.org/10.4171/owr/2018/54}.
\item[S3] Recent simplex lower bounds
\url{https://arxiv.org/abs/2403.04886}.
\textit{Underlying Oberwolfach Reports material retains its source-specific licence, including CC BY-SA 4.0 where applicable; see THIRD\_PARTY\_NOTICES.md.}
\end{itemize}
\paragraph{Preparation and review.}
Original question attributed to the cited source and UnsolvedMath contributors. Codex (GPT-6; exact serving revision not exposed) assisted with definitions, English summaries, source inspection and formatting. No new human mathematical authorship or independent proof verification is asserted. Catalogue identity checked against the complete supplied registry and existing repository statements; known flagged repetitions and obvious identical conjectures excluded. Full mathematical-equivalence/literature review has not been performed. Statement preparation and final review followed contribution prompts C and D; automated checks verify structure and identity, not open status.
\label{end:statement}

\section{Research attempt}
\subsection{Literature, source context and scope}
The full original catalogue body and source qualifications are preserved above. The relevant primary workshop passages were read in the cached publisher report. Historical claims and subsequent attributed results are distinguished from the mathematical work below. A restricted example, standard lemma or known theorem is not counted as new progress on the entire original target. Missing assumptions and repeated or rhetorical extracts are not silently turned into different solved problems.

\subsection{Small graph diameter does not bound every improving pivot path}
A tempting geometric explanation is that a polytope with small graph diameter must force short simplex runs. We test that inference with an explicit deformed cube. Fix $0<\eta<1/2$ and let
\[
P_d=\{x:0\le x_1\le1,\quad \eta x_{i-1}\le x_i\le1-\eta x_{i-1}\ (2\le i\le d)\}.
\]
Maximize $x_d$. The parameter and objective define an auxiliary family, not a distribution of real-world linear programs.
\begin{proposition}
The graph of $P_d$ is the $d$-cube and has diameter $d$, but it has a strictly improving edge path for $x_d$ visiting all $2^d$ vertices.
\end{proposition}
\begin{proof}
Induction on coordinates shows $0\le x_i\le1$. Hence each pair of lower and upper bounds has strictly positive separation, at least $1-2\eta$, and cannot both be tight. For each bit vector $b\in\{0,1\}^d$, select one tight bound in each pair, giving the unique point
\[
x_1=b_1,\qquad x_i=b_i+(1-2b_i)\eta x_{i-1}.
\]
It is feasible and its $d$ active constraints have a triangular, invertible coefficient matrix. It is therefore a vertex. Conversely a vertex must have $d$ independent active constraints; at most one can come from each of the $d$ pairs, so every vertex has this form. Fixing all but one of the selected pairs leaves a one-dimensional face between the two bit choices: the free coordinate ranges between its two distinct bounds, and later fixed coordinates depend affinely on it and stay feasible. Thus bit strings differing in one coordinate are adjacent. Each vertex has exactly $d$ active independent constraints and hence $d$ incident edges, so these are all edges. The graph is the cube, with distances equal to Hamming distances and diameter $d$.
Construct the improving Hamilton path inductively. For $d=1$ traverse 0 then 1. For dimension $d$, first follow the previous path with last bit 0. Along it $x_d=\eta x_{d-1}$ strictly increases from 0 to $\eta$. At its endpoint change the last bit to 1, increasing $x_d$ from $\eta$ to $1-\eta$. Then traverse the previous path in reverse with last bit 1; now $x_d=1-\eta x_{d-1}$ strictly increases to 1. Consecutive bit strings differ in one coordinate, so every step is an edge, and the path visits all vertices exactly once. The middle increase uses $\eta<1/2$.
\end{proof}
\subsection{Verification and remaining gaps}
For $\eta=1/10$ and $1\le d\le8$, all reflected-Gray-code vertices were evaluated with exact rational arithmetic; consecutive objective values strictly increase and consecutive bit labels differ in one position. There are $2^d-1$ pivots despite graph diameter $d$. This disproves only the assertion that small undirected diameter controls every improving rule. It does not preclude a different polynomial pivot rule, does not establish a lower bound for all rules, and does not explain observed practical performance. A rule-specific geometric or probabilistic argument is still needed.

\section{References and source attribution}
Primary workshop report: \url{https://doi.org/10.4171/owr/2018/54}. Relevant passages read: 2018/54 p. 3229, simplex pivot rules, practical behavior and geometric explanation questions.

All retrievals were checked on 11 October 2026. Original question attribution remains in the preserved Sources section. External theorems are attributed and not claimed as new discoveries or fully verified merely by citation.


\paragraph{Integration review.} Codex (GPT-6; exact serving revision not exposed) reviewed the written arguments and expanded layout and source-reference wrappers. Repository metadata and the production web renderer were checked using Python 3.13 and JavaScript with MathJax 3.2.2. This model-assisted review is not independent human proof verification; no PDF compilation or formal proof checking is asserted.

\section{Completion Estimate}
\noindent\textbf{COMPLETION: 0\%}
No substantiated new progress toward the entire original target is claimed. The worked results and failed route are delimited above; a known result, illustrative model, or auxiliary obstruction is not counted as original-target completion.

\section{Final status and attribution}
\textbf{Attempt status: unresolved.} This describes this research attempt, not a new ruling on the source catalogue's status or on external literature claims.
The mathematics and this self-audit were generated by OpenAI \texttt{gpt-6-astra}, reasoning \texttt{ultra}, on 11 October 2026. Codex (GPT-6; exact serving revision not exposed) orchestrated the batch, reviewed the written mathematical arguments, and contributed formatting/repository review. This additional model review is not independent human verification. No human mathematical authorship, independent human proof checking, or formal verification is claimed. The historical preparation paragraphs in the reproduced catalogue body concern that body only.

Contribution prompt A was applied to the complete definition and supplied source record, followed by the adversarial checks of prompt B. Local reproduction files are \texttt{batch022/checks.py}, \texttt{computations.json}, and the per-ID audit JSON. They were not public links at preparation time. Numerical checks are finite and do not prove unrestricted results. No LaTeX engine was found locally; compilation was not performed by this research agent.

\end{document}
